\section{Corruption Experiment: ¿el modelo realmente usa otros datapoints?}

El benchmark performance no demuestra directamente que la attention entre datapoints cumpla una función: un NPT de alta capacidad podría aprender a ignorar las demás rows y degenerar en un per-row parametric model común. Por ello, los autores diseñaron un intervention-like test que, mientras mantiene intacta la target row actual, destruye la información relacional de todas las demás rows y luego observa si la predicción empeora.

\lecturefigure{slide-24-corruption-goal.jpg}{Objetivo del corruption experiment: impedir que el modelo obtenga relaciones útiles de otros datapoints.}{Video oficial de Stanford Online 00:51:16--00:51:31.}

Para cada target entry evaluada $(i,d)$, el método conserva la fila $i$ y aplica a cada attribute column de las demás rows una permutation aleatoria independiente. Puede escribirse como
\[
\widetilde X_{k,j}=X_{\pi_j(k),j},
\qquad k\ne i,
\]
donde cada column $j$ usa una permutation independiente $\pi_j$.

\begin{itemize}
  \item La query row actual $i$ no cambia;
  \item los marginal values de cada column se conservan aproximadamente;
  \item las correspondencias feature--feature y feature--target que existían dentro de una misma row se rompen;
  \item si la predicción depende de la estructura de las demás rows, la performance debería bajar tras la corruption.
\end{itemize}

\lecturefigure{slide-25-corruption-method.jpg}{Corruption method: se permutan las demás rows attribute por attribute y se conserva la query row.}{Video oficial de Stanford Online 00:51:32--00:53:01.}

\teachervoice{El ponente aclaró expresamente que no quería «estropear» la entrada con ruido simple, porque el modelo podría limitarse a detectar un distribution shift. La permutation column por column conserva en lo posible las marginal statistics, pero destruye las correspondencias aprovechables entre campos y entre muestras.}

\subsection{Resultados: algunas tareas dependen fuertemente, otras ignoran activamente}

Si el NPT dependiera por igual de las demás rows en todos los conjuntos de datos, la corruption debería causar pérdidas grandes en general; los resultados reales son más matizados: tareas como Protein empeoran notablemente, mientras que Forest, Kick, Breast Cancer y otras casi no cambian. Esto indica que la architecture ofrece un relation channel, pero la optimization puede decidir, según los datos, si lo usa o no.

\lecturefigure{slide-26-corruption-results.jpg}{Performance change después de destruir la información de los demás datapoints.}{Video oficial de Stanford Online 00:53:02--00:54:43.}

\begin{table}[H]
\centering
\small
\caption{Table 2 del artículo: cambio de performance después de la corruption; los valores negativos indican empeoramiento.}
\begin{tabularx}{\textwidth}{L{0.21\textwidth}*{4}{>{\centering\arraybackslash}X}}
\toprule
Conjunto de datos de clasificación & CIFAR-10 & Poker & Income & Higgs \\
\midrule
$\Delta$ Accuracy & $-1.2$ & $-1.1$ & $-1.1$ & $-0.5$ \\
\bottomrule
\end{tabularx}
\vspace{0.6em}
\begin{tabularx}{\textwidth}{L{0.21\textwidth}*{4}{>{\centering\arraybackslash}X}}
\toprule
Conjunto de datos de clasificación & MNIST & Forest & Kick & Breast Cancer \\
\midrule
$\Delta$ Accuracy & $-0.4$ & $-0.1$ & $-0.1$ & $0.0$ \\
\bottomrule
\end{tabularx}
\vspace{0.6em}
\begin{tabularx}{\textwidth}{L{0.21\textwidth}*{4}{>{\centering\arraybackslash}X}}
\toprule
Conjunto de datos de regresión & Yacht & Protein & Boston & Concrete \\
\midrule
$\Delta\mathrm{RMSE}/\mathrm{RMSE}$ & $-52\%$ & $-21\%$ & $-20\%$ & $-7\%$ \\
\bottomrule
\end{tabularx}
\end{table}

\begin{importantbox}{Lo más interesante no es que «todo baje»}
En algunos datasets, el NPT descubre que las demás rows no aportan suficiente valor, de modo que puede aprender a ignorarlas y degenerar en un predictor aproximadamente parametric; en otros datasets, la corruption provoca un colapso grave de la performance. Lo que se aprende no es la adopción fija de una non-parametric rule, sino el grado de dependencia.
\end{importantbox}

\begin{warningbox}{La corruption respalda la dependencia del mecanismo, pero no es una prueba causal completa}
Un performance drop indica que la predicción aprovechó la structure entre rows que fue destruida; por sí solo, no nos dice cuál es la única ruta causal de cierta attention head, ni permite concluir que «lo non-parametric es en general superior a lo parametric». La corruption también altera la distribución conjunta de orden superior, por lo que todavía podría introducir algún out-of-distribution pattern al que el modelo sea sensible.
\end{warningbox}

\teachervoice{Tanto el artículo como la clase señalan con cautela que el hecho de que Forest, Kick y Breast Cancer casi no se vean afectados no es un fallo del experimento, sino que el modelo posiblemente juzgó que el lookup entre datapoints no ofrecía ventaja. Esta posibilidad de degenerar respalda justamente la afirmación de que «el grado de dependencia se aprende de los datos».}

\subsection{Resumen de la sección}

El corruption test avanza de «el modelo tiene attention entre filas» a «el modelo realmente depende de la información entre filas en varias tareas reales». La evidencia proviene de un performance change de tipo contrafactual, pero el experimento no puede identificar un mecanismo único ni generalizarse a una conclusión universal sobre la superioridad de lo parametric o lo non-parametric.
